% !TeX root = main.tex

Dự án ``Xây dựng nền tảng mạng xã hội rèn luyện thể dục thể thao cộng đồng FitConnect'' đã được hoàn thành với kết quả vượt ngoài mong đợi ban đầu. Nhìn lại toàn bộ hành trình từ ý tưởng đến sản phẩm hoàn chỉnh, đây là một nỗ lực nghiêm túc và có tính tiên phong nhằm lấp đầy khoảng trống mà các nền tảng hiện tại chưa giải quyết được: tích hợp sâu giữa mạng xã hội cộng đồng và các công cụ hỗ trợ rèn luyện sức khỏe cá nhân trong một hệ sinh thái duy nhất, được dẫn dắt bởi Trí tuệ nhân tạo.

\section{Kết quả đạt được}

\subsection{Hiện thực đầy đủ 9 nhóm tính năng cốt lõi}

Hệ thống FitConnect đã được hiện thực thành công toàn bộ 9 nhóm tính năng cốt lõi đã đề ra trong giai đoạn phân tích yêu cầu, mỗi nhóm đều được triển khai đến mức độ sẵn sàng đưa vào sử dụng thực tế.

\textbf{Nhóm 1 -- Quản lý tài khoản và hồ sơ người dùng:} Hệ thống cung cấp đầy đủ quy trình đăng ký, đăng nhập thông qua email lẫn Google OAuth 2.0, cơ chế khôi phục mật khẩu an toàn bằng mã OTP 6 chữ số gửi qua email, và một không gian Profile Dashboard toàn diện cho phép người dùng cá nhân hóa hồ sơ (ảnh đại diện, ảnh nền, thông tin cá nhân, thông số nhân trắc học) đồng thời theo dõi lịch sử toàn bộ hoạt động bao gồm bài viết, sự kiện, thử thách và giáo án tập luyện.

\textbf{Nhóm 2 -- Mạng xã hội cộng đồng:} Bảng tin cộng đồng (Newsfeed) được xây dựng với đầy đủ tính năng của một mạng xã hội hiện đại: đăng bài viết kèm hình ảnh, bày tỏ cảm xúc, bình luận và trả lời bình luận theo luồng, chia sẻ nội dung. Điểm đặc biệt của nhóm tính năng này là sự tích hợp của hai lớp kiểm duyệt nội dung thông minh: mô hình NSFW Detection (TensorFlow.js) chạy trên trình duyệt để phát hiện hình ảnh nhạy cảm trước khi đăng, và bộ lọc từ ngữ tự động để chặn nội dung vi phạm tiêu chuẩn cộng đồng ngay cả trong phần bình luận. Cơ chế báo cáo vi phạm cộng đồng kết hợp với bảng điều khiển kiểm duyệt toàn diện của Quản trị viên tạo nên một vòng kiểm soát chặt chẽ.

\textbf{Nhóm 3 -- Sự kiện thể thao:} Đây là nhóm tính năng phức tạp và đặc trưng nhất của nền tảng. Hệ thống hỗ trợ song song hai hình thức thi đua: Ngoài trời (chạy bộ, đạp xe, leo núi) với tính năng theo dõi lộ trình GPS thời gian thực trên bản đồ Goong Maps, tự động tính toán quãng đường, tốc độ và calo tiêu hao; và Trong nhà (Yoga, Cardio, Gym) với không gian Video Call nhóm tích hợp AI nhận diện khuôn mặt (face-api.js) để tự động giám sát thời gian hiện diện hợp lệ của từng người tham gia, đảm bảo tính công bằng tuyệt đối. Ngoài ra, tính năng AI điền tự động thông tin sự kiện từ một đoạn mô tả ngôn ngữ tự nhiên đã rút ngắn đáng kể thời gian tổ chức.

\textbf{Nhóm 4 -- Thử thách cộng đồng:} Ba loại hình thử thách (Ăn uống/Dinh dưỡng, Thể dục, Ngoài trời) đều được hiện thực với cơ chế check-in tiến độ chuyên biệt. Thử thách dinh dưỡng xác minh ảnh bữa ăn bằng AI đa phương thức; thử thách vận động tích hợp trực tiếp vào không gian tập luyện của nhóm tính năng 6; thử thách ngoài trời sử dụng lại cơ sở hạ tầng GPS đã xây dựng. Hệ thống lịch biểu trực quan giúp theo dõi chuỗi streak liên tục và bảng xếp hạng cộng đồng thời gian thực.

\textbf{Nhóm 5 -- Kết nối mạng lưới bạn bè:} Hệ thống xây dựng cơ chế kết bạn và theo dõi đầy đủ, với bảng tin (News Feed) ưu tiên hiển thị nội dung từ những người trong mạng lưới. Mọi hoạt động thể thao nổi bật như tham gia sự kiện hay chinh phục thử thách đều tự động cập nhật lên bảng tin để lan tỏa động lực trong cộng đồng.

\textbf{Nhóm 6 -- Tập luyện cá nhân:} Hệ thống tập luyện cá nhân cung cấp ba phương thức lựa chọn bài tập bổ sung lẫn nhau: tùy chỉnh thủ công theo thiết bị sẵn có hoặc nhóm cơ mục tiêu (với trực quan hóa nhóm cơ bắp qua thư viện React Body Highlighter), gợi ý AI thông minh dựa trên mô tả tình trạng sức khỏe bằng ngôn ngữ tự nhiên, hoặc sử dụng lại các giáo án đã lưu. Người dùng có thể lên lịch tập luyện hàng tuần và được nhắc nhở tự động, đồng thời tất cả buổi tập được lưu vào lịch sử chi tiết với calo tiêu hao, thời lượng và nhóm cơ đã kích hoạt. Tính năng theo dõi hoạt động thể thao ngoài trời độc lập (chạy bộ, đạp xe) cũng được tích hợp để chuẩn bị cho các sự kiện và thử thách.

\textbf{Nhóm 7 -- Tính toán các chỉ số sức khỏe và gợi ý AI:} Bộ công cụ tính toán y khoa bao gồm 8 chỉ số: BMI, BMR, TDEE, Body Fat, IBW (Cân nặng lý tưởng), LBM (Khối lượng cơ thể nạc), Calo đốt cháy và Lượng nước cần mỗi ngày. Mỗi lần tính toán đều được lưu vào lịch sử và hiển thị qua biểu đồ xu hướng theo thời gian. Đặc biệt, tính năng AI Phân Tích Sức Khỏe tổng hợp toàn bộ chỉ số để đưa ra lời khuyên dinh dưỡng và danh sách bài tập gợi ý cá nhân hóa từ kho dữ liệu của hệ thống.

\textbf{Nhóm 8 -- Lịch trình cá nhân và thông báo:} Giao diện Lịch cá nhân (Calendar) tự động đồng bộ tất cả hoạt động mà người dùng đã đăng ký -- sự kiện thể thao, thử thách cộng đồng, lịch tập luyện cá nhân -- vào một giao diện thống nhất. Hệ thống thông báo thời gian thực (Socket.io) gửi cảnh báo khi đến giờ tập, khi bạn bè tương tác và khi có mục tiêu sắp đến hạn.

\textbf{Nhóm 9 -- Quản trị hệ thống toàn diện:} Bảng điều khiển Quản trị viên (Admin Dashboard) tích hợp biểu đồ thống kê thời gian thực về tăng trưởng người dùng, mức độ tương tác theo từng nhóm tính năng và thống kê sử dụng AI theo danh mục. Quản trị viên có toàn quyền kiểm duyệt và xử lý mọi nội dung vi phạm (bài viết, sự kiện, thử thách, tài khoản) thông qua luồng báo cáo từ cộng đồng, đồng thời quản lý toàn bộ tài nguyên dữ liệu nền tảng bao gồm danh mục thể thao, thiết bị tập luyện và kho bài tập.

\subsection{Kết quả kiểm thử}

Chất lượng của hệ thống được kiểm chứng thông qua nhiều phương pháp kiểm thử đa cấp độ:

\begin{itemize}
    \item \textbf{Kiểm thử chức năng:} 65 kịch bản kiểm thử UAT (User Acceptance Testing) bao phủ toàn bộ 9 nhóm tính năng đều cho kết quả \textbf{96.9\% thành công} (63/65 kịch bản passed). 15 kịch bản Happy Case đại diện (TC-01 đến TC-15) đều vượt qua hoàn toàn.

    \item \textbf{Kiểm thử Unit và API:} Tổng cộng 80 Unit Test với tỉ lệ bao phủ mã nguồn trung bình \textbf{88.5\%}; 73 API Endpoint được kiểm thử, đạt tỉ lệ thành công \textbf{97.3\%} (71/73 endpoints passed).

    \item \textbf{Kiểm thử tải (Loader.io):} Hệ thống xử lý ổn định \textbf{1.000 người dùng đồng thời} trong vòng 1 phút với thời gian phản hồi trung bình chỉ \textbf{13ms}, tỉ lệ lỗi \textbf{0\%} (953/953 yêu cầu thành công), chứng minh kiến trúc Event Loop của Node.js hoạt động hiệu quả dưới áp lực cao.

    \item \textbf{Chất lượng giao diện (Google PageSpeed Insights):} Frontend đạt điểm số xuất sắc -- \textbf{Hiệu suất 85/100}, \textbf{Khả năng tiếp cận 92/100}, \textbf{Best Practices 100/100} và \textbf{SEO 100/100} -- đáp ứng toàn bộ các tiêu chuẩn kỹ thuật web hiện đại năm 2026.

    \item \textbf{Hiệu năng cơ sở dữ liệu (MongoDB Atlas):} Cụm Replica Set 3 nodes vận hành ổn định liên tục trong 2 tháng (tháng 3 đến tháng 4 năm 2026). Thời gian phản hồi các truy vấn được tối ưu rõ rệt: truy vấn danh sách bài viết phân trang đạt \textbf{$\sim$80ms}, truy vấn tổng hợp tiến độ Thử thách (Pipeline Aggregation) đạt \textbf{$\sim$150ms}, tìm kiếm bài tập thể dục (Text Indexing) đạt \textbf{$\sim$45ms}.

    \item \textbf{Kiểm thử bảo mật:} Hệ thống vượt qua tất cả kịch bản tấn công: từ chối Token hết hạn/bị giả mạo (401 Unauthorized), ngăn chặn XSS và NoSQL Injection, phân quyền phân đoạn chặt chẽ (User không thể truy cập route \texttt{/api/admin}), và cơ chế Rate Limiting tự động chặn IP khi vượt 100 request/15 phút.
\end{itemize}

\subsection{Thành tựu về mặt kỹ thuật}

Về mặt kỹ thuật, hệ thống được xây dựng trên kiến trúc hiện đại và đồng nhất ngôn ngữ từ đầu đến cuối bằng \textbf{JavaScript/TypeScript}. Backend (Node.js + Express.js) triển khai đúng mô hình Routes~$\rightarrow$~Controllers~$\rightarrow$~Services ba lớp với MongoDB Atlas (Replica Set) làm cơ sở dữ liệu chính; Frontend người dùng (React 18 + Vite + TanStack Query) và Frontend quản trị viên (React 18) được triển khai độc lập theo mô hình monorepo. Trí tuệ nhân tạo được tích hợp xuyên suốt ở nhiều điểm chiến lược: kiểm duyệt hình ảnh (NSFW Detection), nhận diện khuôn mặt (face-api.js), xác minh bữa ăn (AI đa phương thức), phân tích sức khỏe và gợi ý cá nhân hóa, tạo nội dung tự động cho sự kiện và thử thách. Toàn bộ hệ thống AI sử dụng mô hình Proxy Pattern để che giấu API Key phía Backend và ghi nhận thống kê sử dụng theo từng danh mục tính năng.

\section{Hạn chế hiện tại}

Bên cạnh những kết quả đạt được, dự án còn tồn tại một số hạn chế khách quan xuất phát từ giới hạn về thời gian và nguồn lực của một đồ án tốt nghiệp:

\begin{itemize}
    \item \textbf{Quy mô cơ sở dữ liệu bài tập:} Kho dữ liệu bài tập hiện tại được xây dựng thủ công với số lượng bài tập còn hạn chế, chưa bao phủ đầy đủ tất cả các môn thể thao và phương pháp tập luyện chuyên sâu. Điều này ảnh hưởng đến chất lượng gợi ý AI trong một số trường hợp đặc thù.

    \item \textbf{Tính cá nhân hóa AI còn ở mức tổng quát:} Các gợi ý bài tập và phân tích sức khỏe hiện tại sử dụng mô hình ngôn ngữ lớn (LLM) kết hợp dữ liệu người dùng, nhưng chưa có cơ chế học tăng cường (Reinforcement Learning) để AI cải thiện dần độ chính xác gợi ý dựa trên phản hồi thực tế từ hành vi tập luyện lâu dài của người dùng.

    \item \textbf{Chưa có ứng dụng di động riêng:} Hệ thống hiện chỉ là ứng dụng web với thiết kế Responsive, chưa có ứng dụng iOS/Android native. Các tính năng đòi hỏi quyền truy cập phần cứng sâu hơn (cảm biến gia tốc, GPS background tracking, thông báo đẩy) còn bị giới hạn bởi API trình duyệt.

    \item \textbf{Năng lực xử lý Video Call ở quy mô lớn:} Hạ tầng WebRTC hiện tại được triển khai theo mô hình Peer-to-Peer (P2P) thông qua máy chủ TURN. Khi số lượng người tham gia một phòng Video Call vượt ngưỡng, độ trễ và chất lượng hình ảnh có thể bị ảnh hưởng; giải pháp SFU (Selective Forwarding Unit) chuyên dụng chưa được tích hợp.

    \item \textbf{Mức độ tự động hóa kiểm thử:} Bộ kiểm thử hiện tại chủ yếu là Unit Test và kiểm thử API thủ công; chưa có bộ kiểm thử End-to-End (E2E) tự động bằng Playwright/Cypress để phủ toàn bộ các luồng người dùng phức tạp.
\end{itemize}

\section{Định hướng phát triển tương lai}

Với nền tảng kỹ thuật vững chắc và kiến trúc mở rộng được đã xây dựng, FitConnect sở hữu tiềm năng phát triển rõ ràng theo nhiều hướng chiến lược:

\begin{itemize}
    \item \textbf{Phát triển ứng dụng di động (Mobile App):} Xây dựng ứng dụng iOS và Android bằng React Native hoặc Flutter để tận dụng đầy đủ năng lực phần cứng của thiết bị thông minh: GPS Background Tracking (theo dõi liên tục khi tắt màn hình), kết nối Bluetooth với thiết bị đo nhịp tim và cảm biến, thông báo đẩy (Push Notification) theo thời gian thực, và quét mã vạch thực phẩm để tra cứu thông tin dinh dưỡng tức thì.

    \item \textbf{Nâng cấp hạ tầng Video Call (SFU):} Tích hợp máy chủ SFU chuyên dụng (mediasoup hoặc LiveKit) để thay thế mô hình P2P hiện tại, cho phép phòng Video Call hỗ trợ hàng trăm người tham gia đồng thời với chất lượng hình ảnh ổn định, cùng các tính năng nâng cao như ghi hình phiên tập và phát lại sau buổi sự kiện.

    \item \textbf{AI cá nhân hóa thế hệ tiếp theo:} Xây dựng mô hình học máy (Machine Learning) dựa trên lịch sử tập luyện dài hạn của từng người dùng để AI dự đoán được nguy cơ chấn thương, tự động điều chỉnh khối lượng tập theo chu kỳ phục hồi (Periodization), và đề xuất lộ trình dài hạn (12 tuần, 6 tháng) tùy theo mục tiêu cụ thể.

    \item \textbf{Tích hợp thiết bị đeo thông minh (Wearables):} Kết nối API với Apple Health, Google Fit và các thiết bị smartwatch (Apple Watch, Garmin, Fitbit) để tự động đồng bộ dữ liệu nhịp tim, bước chân, chất lượng giấc ngủ và calo tiêu hao thực tế, làm phong phú thêm hồ sơ sức khỏe và tăng độ chính xác của gợi ý AI.

    \item \textbf{Mở rộng tính năng Dinh dưỡng:} Xây dựng cơ sở dữ liệu dinh dưỡng món ăn Việt Nam hợp tác với Viện Dinh dưỡng Quốc gia, tích hợp tính năng nhận diện món ăn từ hình ảnh (Food Recognition AI) và lên kế hoạch thực đơn cá nhân hóa hoàn chỉnh kết hợp với chỉ tiêu TDEE đã tính.

    \item \textbf{Gamification và hệ thống phần thưởng:} Phát triển hệ thống điểm danh tiếng (Reputation Points), huy hiệu thành tích (Achievement Badges) và bảng xếp hạng toàn cầu theo từng môn thể thao để tăng động lực dài hạn. Tích hợp cơ chế thưởng cộng đồng (Community Rewards) và các chương trình đối tác với thương hiệu thể thao.

    \item \textbf{Nền tảng cho Huấn luyện viên chuyên nghiệp (Coach Platform):} Xây dựng vai trò Huấn luyện viên được xác minh (Verified Coach) với khả năng tạo và bán giáo án tập luyện chuyên nghiệp, tổ chức các lớp học trực tuyến có phí thông qua Video Call, và theo dõi tiến độ học viên theo thời gian thực.
\end{itemize}

\section{Kết luận}

Nhìn lại toàn bộ quá trình thực hiện, đồ án đã hoàn thành sứ mệnh của mình: xây dựng một nền tảng mạng xã hội sức khỏe cộng đồng toàn diện mang tên FitConnect, vừa đáp ứng đầy đủ 9 nhóm tính năng đã cam kết, vừa được kiểm chứng chất lượng nghiêm túc qua nhiều phương pháp kiểm thử khoa học. Hệ thống không chỉ là một sản phẩm kỹ thuật đơn thuần mà còn là một giải pháp số có ý nghĩa xã hội thực tiễn: giúp người Việt Nam dễ dàng hơn trong việc bắt đầu, duy trì và lan tỏa thói quen rèn luyện thể thao lành mạnh thông qua sức mạnh của cộng đồng và công nghệ AI.

Với kiến trúc phần mềm hiện đại, mã nguồn được tổ chức sạch sẽ theo chuẩn TypeScript, và cơ sở hạ tầng cơ sở dữ liệu phân tán đã được kiểm chứng, FitConnect có đầy đủ nền tảng kỹ thuật vững chắc để tiếp tục phát triển theo những định hướng đã vạch ra, hướng tới trở thành một giải pháp sức khỏe cộng đồng toàn diện đồng hành cùng người Việt Nam trên hành trình xây dựng cuộc sống khỏe mạnh và chất lượng hơn.
